%%%%%%%%%%%%%%%%%%%%%%%%%%%%%%%%%%%%%%%%%%%%%%%%%%%%%%%%%%%%%%%%%%%%%%%%%%%%%%%
%  WCNPS 2026 - 11th Workshop on Communication Networks and Power Systems
%
%  Migracao Pos-Quantica em Pipelines DevSecOps: o Custo do Handshake mTLS
%  e um Degrau na Janela de Congestionamento
%%%%%%%%%%%%%%%%%%%%%%%%%%%%%%%%%%%%%%%%%%%%%%%%%%%%%%%%%%%%%%%%%%%%%%%%%%%%%%%

\documentclass[conference]{IEEEtran}
\IEEEoverridecommandlockouts
\renewcommand\IEEEkeywordsname{Palavras-chave}
\usepackage[utf8]{inputenc}
\usepackage[T1]{fontenc}
\usepackage[american,brazil]{babel}
\usepackage[center]{caption}
\usepackage{graphicx}
\usepackage{amsmath,amssymb}
\usepackage{booktabs}
\usepackage{scalerel}
\usepackage{tikz}
\usetikzlibrary{svg.path}

% ------------------------------------------------------------------ ORCID
\definecolor{orcidlogocol}{HTML}{A6CE39}
\tikzset{
  orcidlogo/.pic={
    \fill[orcidlogocol] svg{M256,128c0,70.7-57.3,128-128,128C57.3,256,0,198.7,0,128C0,57.3,57.3,0,128,0C198.7,0,256,57.3,256,128z};
    \fill[white] svg{M86.3,186.2H70.9V79.1h15.4v48.4V186.2z}
                 svg{M108.9,79.1h41.6c39.6,0,57,28.3,57,53.6c0,27.5-21.5,53.6-56.8,53.6h-41.8V79.1z M124.3,172.4h24.5c34.9,0,42.9-26.5,42.9-39.7c0-21.5-13.7-39.7-43.7-39.7h-23.7V172.4z}
                 svg{M88.7,56.8c0,5.5-4.5,10.1-10.1,10.1c-5.6,0-10.1-4.6-10.1-10.1c0-5.6,4.5-10.1,10.1-10.1C84.2,46.7,88.7,51.3,88.7,56.8z};
  }
}
\newcommand\orcid[1]{\href{https://orcid.org/#1}{\mbox{\scalerel*{
\begin{tikzpicture}[yscale=-1,transform shape]
\pic{orcidlogo};
\end{tikzpicture}
}{|}}}}

\usepackage[bookmarks=false]{hyperref}

% ------------------------------------------------- cabecalho do workshop
\makeatletter
\let\old@ps@headings\ps@headings
\let\old@ps@IEEEtitlepagestyle\ps@IEEEtitlepagestyle
\def\confheader#1{%
\def\ps@headings{\old@ps@headings%
\def\@oddhead{\strut\hfill#1\hfill\strut}%
\def\@evenhead{\strut\hfill#1\hfill\strut}}%
\def\ps@IEEEtitlepagestyle{\old@ps@IEEEtitlepagestyle%
\def\@oddhead{\strut\hfill#1\hfill\strut}%
\def\@evenhead{\strut\hfill#1\hfill\strut}}%
\ps@headings}
\makeatother
\confheader{11th Workshop on Communication Networks and Power Systems (WCNPS 2026)}

\begin{document}

\title{Migração Pós-Quântica em Pipelines DevSecOps: o Custo do
\textit{Handshake} mTLS é um Degrau na Janela de Congestionamento}

\author{\IEEEauthorblockN{Jaime de Melo Gama da Silva \orcid{0009-0007-8925-4056}\IEEEauthorrefmark{1}\IEEEauthorrefmark{2},
Daniel Chaves Café \orcid{0000-0003-1132-3430}\IEEEauthorrefmark{2},
Rômulo F. dos Santos \orcid{0009-0001-3588-7583}\IEEEauthorrefmark{1}}
\IEEEauthorblockA{Universidade de Brasília, Faculdade de Tecnologia -- Departamento de Engenharia Elétrica\IEEEauthorrefmark{2}\\
Email: jaimemelo671@gmail.com, dcafe@unb.br, romulodba@gmail.com\IEEEauthorrefmark{1}}
}

\IEEEpubid{\makebox[\columnwidth]{979-8-3315-7075-0/25/\$31.00 $\copyright$2025 IEEE} \hspace{\columnsep}\makebox[\columnwidth]{ }}

\maketitle

% =============================================================== RESUMO
\begin{abstract}
Pipelines DevSecOps substituíram segredos estáticos por identidades de carga de
trabalho de vida curta, o que faz cada tarefa pagar um \textit{handshake} mTLS
completo, a frio e com cadeia de certificados dos dois lados. A literatura de
desempenho de criptografia pós-quântica em TLS foi construída sobre o perfil da
Web, com autenticação unilateral e sessões longas, e não cobre esse regime. Este
artigo caracteriza o custo da migração num fatorial $2\times2$ de troca de chaves
e autenticação, com 30 execuções independentes por configuração e
\textit{bootstrap} hierárquico, sobre plataforma nativa deliberadamente modesta.
O orçamento de bytes é aditivo, com resíduo de 6 bytes em 21\,339, e concorda com
a contagem normativa de campos dentro de sete bytes, o que permite prever o
volume de combinações não medidas. Sob 40~ms de ida e volta com 0,5\% de perda,
porém, a latência deixa de acompanhar o tamanho e vira um degrau: a configuração
cujos voos excedem a janela inicial de congestionamento leva 84\% mais tempo que
a linha de base clássica, contra 8\% da imediatamente abaixo do limiar. Uma
varredura da janela revela uma escada de dois degraus iguais, consequência de a
autenticação mútua aplicar o critério a cada sentido, e permite estimar o tamanho
de segmento efetivo a partir da própria latência. Como a plataforma infla o termo
criptográfico, o efeito é um limite inferior para hardware de maior porte.
\end{abstract}

\begin{IEEEkeywords}
DevSecOps, criptografia pós-quântica, TLS, controle de congestionamento,
ML-KEM, ML-DSA
\end{IEEEkeywords}

% =============================================================== INTRODUCAO
\section{Introdução}

A automação de construção, teste e implantação de software acumulou um volume de
privilégio que raramente aparece nos modelos de ameaça das organizações. Um
\emph{runner} de CI/CD publica artefatos, altera \textit{clusters} de produção e
assume papéis amplos em nuvem, historicamente apoiado em segredos estáticos que
vazam com frequência em repositórios públicos~\cite{meli2019git} e cuja
exploração consolidou a cadeia de suprimentos como vetor recorrente de
ataque~\cite{ohm2020backstabber}. A resposta consolidada foi substituir esses
segredos por identidades de carga de trabalho emitidas em tempo de execução
mediante atestação, no espírito do SPIFFE e da arquitetura \emph{Zero
Trust}~\cite{rose2020zta,avirneni2025workload}.

Essa mudança tem uma consequência de rede que costuma passar despercebida. Como a
credencial vale segundos e é emitida por tarefa, e como a comunicação entre
componentes do pipeline é mutuamente autenticada, cada conexão paga um
\textit{handshake} completo, a frio, sem reuso de sessão e com cadeia de
certificados apresentada pelos \emph{dois} lados. Numa construção que abre
centenas de conexões curtas, o custo do estabelecimento deixa de ser amortizável
e passa a ser proporcional ao número de tarefas; em malhas de serviço com
\textit{sidecars}~\cite{ashok2021servicemesh} ele é pago a cada salto.

Sobre esse cenário incide a padronização do ML-KEM~\cite{nist2024fips203}, do
ML-DSA~\cite{nist2024fips204} e do SLH-DSA~\cite{nist2024fips205}, que
transformou a migração pós-quântica em projeto de
engenharia~\cite{shor1997polynomial,mosca2018cybersecurity}. A literatura de
desempenho de PQC em TLS é
sólida~\cite{sikeridis2020postquantum,paquin2020benchmarking,schwabe2020kemtls},
porém construída sobre o perfil da Web: autenticação unilateral, sessões longas e
certificados com meses de validade. O perfil de um pipeline é quase o oposto, e a
diferença não é só de escala. Com autenticação mútua e cadeia de dois níveis, o
volume do estabelecimento aproximadamente dobra, o que aproxima o
\textit{handshake} de limiares do transporte que no perfil Web ficam distantes.

Daí a pergunta que orienta o trabalho: \emph{nesse regime, quanto custa a
migração pós-quântica, e onde o custo aparece?} As contribuições são três.
Mostra-se que troca de chaves e autenticação compõem seus custos de bytes por
soma, com resíduo abaixo de 0,03\%, e que o valor absoluto é recuperável dos
tamanhos normativos dos esquemas, o que fornece um modelo preditivo. Demonstra-se
que em enlace com atraso e perda o custo de latência não é proporcional ao
tamanho, mas um degrau determinado pela janela inicial de congestionamento,
verificado por varredura controlada desse parâmetro e não inferido por
correlação. E dos dois anteriores deriva-se uma regra de projeto de cadeia
aplicável antes de qualquer medição.

% =========================================================== REFERENCIAL
\section{Trabalhos Relacionados}
\label{sec:referencial}

Três resultados delimitam o problema. Paquin \textit{et
al.}~\cite{paquin2020benchmarking} mostram que, sob perda de pacotes, o tamanho
das mensagens domina o custo computacional dos esquemas pós-quânticos, e é dessa
observação que parte a suspeita de que enlaces ideais produzem conclusões
otimistas. Sikeridis \textit{et al.}~\cite{sikeridis2020postquantum} demonstram
que a autenticação, e não a troca de chaves, é o componente que mais penaliza o
\textit{handshake} em TLS~1.3, mas o fazem com autenticação unilateral em cenário
Web. O KEMTLS~\cite{schwabe2020kemtls} evidencia que substituir assinaturas por
encapsulamento de chaves reduz banda e ciclos. Trabalhos sobre certificados X.509
pós-quânticos~\cite{kampanakis2018viability} e sobre transição de infraestrutura
de chave pública~\cite{bindel2017transitioning} já antecipavam o crescimento das
cadeias como obstáculo dominante e a não linearidade do efeito, mas nenhum
identifica o mecanismo de transporte responsável nem o parâmetro que o controla.

Há ainda uma tradição de aferição de PQC em plataformas restritas cujo protocolo
metodológico este trabalho adota. O pqm4~\cite{kannwischer2019pqm4} afere os
candidatos do NIST num microcontrolador Cortex-M4 de 196~KiB, e
Bürstinghaus-Steinbach \textit{et al.}~\cite{burstinghaus2020embedded} integram
Kyber e SPHINCS+ ao mbed~TLS sobre quatro plataformas embarcadas. Em ambos, a
plataforma modesta não é limitação declarada com humildade, e sim instrumento que
estabelece viabilidade por cima.

Do lado da automação, revisões de DevSecOps apontam a automação como o desafio
mais relatado~\cite{rajapakse2022challenges}, e a literatura de identidade de
carga de trabalho~\cite{avirneni2025workload} trata PQC de passagem. Na revisão
realizada não se encontrou caracterização do custo de PQC em autenticação mútua
com cadeias de dois níveis, regime praticado por pipelines e malhas de serviço.

% ============================================================ BANCADA
\section{Bancada e Desenho Experimental}
\label{sec:bancada}

\subsection{Escolha da plataforma}
\label{sec:plataforma}

A campanha foi executada sobre um Intel Core i3-4005U de dois núcleos a 1,70~GHz,
microarquitetura de 2013, sem Turbo Boost e com 8~GB de DDR3. A escolha é
metodológica e apoia-se em três argumentos.

O primeiro é o de limite superior, como em pqm4~\cite{kannwischer2019pqm4} e em
Bürstinghaus-Steinbach \textit{et al.}~\cite{burstinghaus2020embedded}: um custo
aceitável na plataforma mais lenta é aceitável em qualquer plataforma mais
rápida, e a recíproca não vale. A CPU empregada situa-se entre os extremos
usualmente reportados, cerca de duas ordens de grandeza acima de um
microcontrolador e uma ordem abaixo de um nó de servidor contemporâneo.

O segundo é específico deste trabalho e mais forte. O resultado central é que, em
enlace com atraso e perda, o custo passa a ser dominado por um termo de rede e
não pelo criptográfico. Uma CPU lenta \emph{infla} o termo criptográfico e,
portanto, \emph{dificulta} observar o efeito. Aqui a criptografia responde por
4,5 a 7,7~ms do \textit{handshake} contra dois tempos de ida e volta de cerca de
40~ms cada, razão em torno de dez para um; num servidor com implementações
vetorizadas o termo criptográfico cairia para fração de milissegundo e a razão
passaria de cem para um. A plataforma é o caso conservador para a tese, e as
magnitudes reportadas são limite inferior do efeito.

O terceiro é de controle experimental. A ausência de Turbo Boost elimina a
variação de frequência que, em processadores de servidor, se confunde com a
variabilidade do algoritmo sob medição; o chassi de baixa potência permite
verificar por contador que não houve \textit{throttling}; e a máquina dedicada
permite isolar um núcleo sem disputa com outros inquilinos, o que em nó
compartilhado raramente é possível. Como virtualização e contêineres degradam
sobretudo a cauda da distribuição~\cite{felter2015updated}, a plataforma nativa
compra precisão onde ela é mais escassa.

\subsection{Protocolo de medição}
\label{sec:protocolo}

Quatro elementos do protocolo daqueles trabalhos foram adotados, e o quinto é
exigido pelo objeto.

\textbf{Frequência fixa e relato em ciclos.} O pqm4 opera a placa em frequência
fixa e reporta ciclos, não tempo de parede, porque ciclos comparam-se entre
plataformas de relógios distintos. Aqui a frequência é constante por projeto, de
modo que ciclos e tempo relacionam-se por 1\,700 ciclos por microssegundo e os
tempos podem ser confrontados com aferições de outras plataformas. Como
calibração, a assinatura isolada com ML-DSA-44 custa 1,95~Mciclos nessa CPU,
contra 0,11~Mciclos do ECDSA P-256.

\textbf{Integração em biblioteca real.} Como em Bürstinghaus-Steinbach \textit{et
al.}, a medição principal é do \textit{handshake} completo sobre a pilha TLS de
produção, e não da soma de primitivas isoladas. Isso captura custos de
serialização e de transporte que a soma ignora, e é neles que reside o efeito
relatado.

\textbf{Linha de base clássica em toda comparação.} Todo resultado é razão em
relação a C1, medida na mesma sessão e na mesma ordem randomizada, sem comparação
contra número de literatura obtido em outra plataforma.

\textbf{Arcabouço reproduzível.} A campanha é disparada por uma linha, com
semente de randomização fixa registrada.

\textbf{Voo por direção, além do total.} Como a hipótese envolve um limiar de
transporte, o total trocado é insuficiente: é o voo de um único lado que precisa
caber na primeira janela.

\subsection{Configuração e instrumentação}
\label{sec:config}

O sistema é o Rocky Linux 10.2 em modo FIPS com OpenSSL 3.5.5, cujo provedor
padrão traz ML-KEM e ML-DSA nativos. Aplicaram-se governador em
\texttt{performance}, \textit{hyper-threading} e estados C profundos desligados,
isolamento do núcleo de medição, ASLR desabilitado, \textit{swap} desligado e
páginas transparentes desabilitadas, medidas que atacam as fontes de viés
documentadas por Mytkowicz \textit{et al.}~\cite{mytkowicz2009wrong}. O relógio é
o \texttt{CLOCK\_MONOTONIC\_RAW}.

Cliente e servidor ocupam \textit{network namespaces} e núcleos distintos,
ligados por par \texttt{veth} com MTU de 1500 bytes e \textit{offloads}
desabilitados: com TSO, GSO e GRO ligados, a agregação de segmentos faria a
contagem de bytes no \textit{socket} deixar de refletir o que a rede veria. O
controle de congestionamento é o CUBIC~\cite{ha2008cubic}, padrão do Linux, e a
opção de \textit{timestamps} está habilitada. Registre-se que o resultado da
Seção~\ref{sec:degrau} não depende da disciplina escolhida: um \textit{handshake}
de poucas dezenas de kilobytes termina ainda dentro da partida lenta, regime em
que CUBIC, Reno e BBR partem da mesma janela inicial e a duplicam, divergindo
apenas na prevenção de congestionamento, à qual esse tamanho não chega. Dois
perfis de enlace foram avaliados: o \texttt{lan} usa o par \texttt{veth} sem
emulação, e o \texttt{wan} aplica \texttt{netem} nos dois sentidos, com 20~ms de
atraso e desvio de 2~ms por sentido e 0,5\% de perda, alvo de 40~ms de ida e
volta, medido entre 38,0 e 42,4~ms conforme a sessão. O \textit{script} de rede
afere esse valor por \textit{ping} após aplicar o perfil e recusa-se a prosseguir
quando o enlace não corresponde ao pretendido, salvaguarda adotada depois de uma
campanha inteira ter sido coletada com o perfil rotulado e não aplicado.

\begin{table}[!t]
\caption{Configurações avaliadas. As quatro primeiras formam o fatorial
$2\times2$: as linhas variam o KEM sobre a mesma curva e as colunas variam a
assinatura. Categoria de segurança NIST entre parênteses.}
\label{tab:tratamentos}
\centering
\footnotesize
\setlength{\tabcolsep}{4pt}
\begin{tabular}{@{}llll@{}}
\toprule
\textbf{ID} & \textbf{Grupo TLS} & \textbf{Assinatura e CA} & \textbf{Papel} \\
\midrule
C1 & \texttt{secp256r1}         & ECDSA P-256 (clás.) & linha de base \\
C3 & \texttt{secp256r1}         & ML-DSA-44 (2)       & fatorial \\
H1 & \texttt{SecP256r1MLKEM768} & ECDSA P-256 (clás.) & fatorial \\
H2 & \texttt{SecP256r1MLKEM768} & ML-DSA-44 (2)       & fatorial \\
P3 & \texttt{MLKEM768}          & ML-DSA-65 (3)       & pós-quântica pura \\
\bottomrule
\end{tabular}
\end{table}

A Tabela~\ref{tab:tratamentos} lista as configurações, com os nomes de grupo tal
como passados a \texttt{SSL\_CTX\_set1\_groups\_list}. Todas usam TLS~1.3 com
autenticação mútua e cadeia homogênea de dois níveis, com a autoridade assinando
a folha com o mesmo esquema desta, de modo que em P3 nem a assinatura da
autoridade nem o \texttt{CertificateVerify} recorrem a esquema clássico. O modo
FIPS retira o X25519 dos grupos disponíveis e obriga a usar \texttt{secp256r1}, o
que tornou o desenho mais limpo do que o planejado: C1 e H1 passam a diferir
\emph{somente} pelo acréscimo do ML-KEM sobre a mesma curva, e o mesmo vale para
C3 e H2.

Os certificados de entidade final carregam, além dos campos usuais, uma extensão
crítica de 134 bytes específica da implantação, idêntica em todas as
configurações. Por ser constante ela não afeta nenhuma comparação, mas integra os
voos reportados e convém registrá-la, porque o critério da
Seção~\ref{sec:degrau} é sensível a essa ordem de grandeza: sem ela o voo do
servidor em P3 mediria 15\,908 bytes e caberia em onze segmentos, e não em doze.
Cadeias de produção costumam carregar extensões próprias de porte semelhante.

Foram medidas duas variáveis: $T_{\mathrm{hs}}$, o \textit{handshake} mTLS a frio
até o primeiro byte de aplicação, e $B_{\mathrm{hs}}$, os bytes nos contadores do
BIO de \textit{socket}, isto é, a carga útil TCP nos dois sentidos, incluindo
cabeçalhos de registro TLS e excluindo cabeçalhos TCP/IP e retransmissões. Cada
conexão paga um \textit{handshake} completo, sem reuso de sessão nem
\textit{tickets}. Cada configuração recebe 30 execuções independentes em
\texttt{lan} e 15 em \texttt{wan}, com 300 \textit{handshakes} e 20 de aquecimento
por execução, em ordem randomizada com semente fixa; o contador de
\textit{throttling} é lido antes e depois de cada sessão, que é descartada se
mudar. A incerteza vem de \textit{bootstrap} hierárquico no formato de Kalibera e
Jones~\cite{kalibera2013rigorous}, reamostrando execuções no primeiro nível e
operações no segundo, com 10\,000 reamostragens: tratar as 9\,000 operações de uma
configuração como independentes produziria intervalos artificialmente estreitos,
já que operações do mesmo lote compartilham estado de \textit{cache} e de
alocador.

% ============================================================ RESULTADOS
\section{O Orçamento de Bytes é Aditivo}
\label{sec:bytes}

\begin{table}[!t]
\caption{Latência e volume do \textit{handshake}, razões em relação a C1 entre
parênteses. Todos os intervalos de 95\% ficam a menos de 1\% da mediana. A última
coluna é o voo do servidor, isto é, os bytes que o cliente recebe.}
\label{tab:resultados}
\centering
\footnotesize
\setlength{\tabcolsep}{3.5pt}
\begin{tabular}{@{}lrrrr@{}}
\toprule
\textbf{ID} &
\textbf{$T_{\mathrm{hs}}$ \texttt{lan}} &
\textbf{$T_{\mathrm{hs}}$ \texttt{wan}} &
\textbf{$B_{\mathrm{hs}}$} &
\textbf{voo serv.} \\
 & ($\mu$s) & ($\mu$s) & (bytes) & (bytes) \\
\midrule
C1 & 4\,516 (1{,}00) & 85\,455 (1{,}00) & 2\,821 (1{,}00)  & 1\,536 \\
C3 & 7\,101 (1{,}57) & 91\,742 (1{,}07) & 21\,890 (7{,}76) & 11\,070 \\
H1 & 5\,149 (1{,}14) & 86\,262 (1{,}01) & 5\,085 (1{,}80)  & 2\,624 \\
H2 & 7\,659 (1{,}70) & 92\,567 (1{,}08) & 24\,160 (8{,}56) & 12\,160 \\
\midrule
P3 & 7\,703 (1{,}71) & 157\,656 (1{,}84) & 31\,918 (11{,}31) & \textbf{16\,042} \\
\bottomrule
\end{tabular}
\end{table}

A Tabela~\ref{tab:resultados} consolida as medições. Um indicador de validade
interna precede a leitura: em C3, H2 e P3 as 9\,000 medições de $B_{\mathrm{hs}}$
de cada configuração produziram um \emph{único} valor distinto, e em C1 e H1
apareceram seis valores num intervalo de cinco bytes, que é a assinatura da
codificação DER do ECDSA, cujos inteiros $r$ e $s$ perdem zeros à esquerda ou
ganham um octeto de preenchimento conforme o bit mais significativo. ML-DSA tem
assinatura de comprimento fixo e não exibe a oscilação.

Com as quatro células do fatorial, os efeitos principais são estimáveis sem
depender da ordem de composição. Em bytes, acrescentar o ML-KEM sobre a mesma
curva custa 2\,270 e trocar ECDSA por ML-DSA-44 custa 19\,075, o que atribui
89,4\% do acréscimo de C1 a H2 à autenticação e 10,6\% à troca de chaves,
confirmando o achado de Sikeridis \textit{et al.}~\cite{sikeridis2020postquantum}
agora em autenticação mútua.

O resultado mais útil é que os dois fatores se compõem por soma. Somando à linha
de base os efeitos medidos isoladamente, prevê-se para H2 um total de 24\,154
bytes contra 24\,160 observados, resíduo de 6 bytes em 21\,339 de acréscimo, da
mesma ordem da oscilação DER acima. Não há interação mensurável na escala em que
os fatores operam, e a aparente interação que surge ao expressar os efeitos como
razões é artefato dessa parametrização.

Os dois efeitos reproduzem, além disso, aritmética de tamanho de campo. Os
2\,270 bytes da troca de chaves correspondem à soma da chave pública do
ML-KEM-768, de 1\,184 bytes, com o texto cifrado, de 1\,088. Os 19\,075 da
autenticação correspondem, dentro de sete bytes, à contagem dos quatro
certificados e dos dois \texttt{CertificateVerify} de um \textit{handshake} mútuo
com cadeia de dois níveis, em que cada chave pública passa de 65 para 1\,312
bytes e cada assinatura de cerca de 71 para 2\,420. Disso resulta um modelo
simples: para cadeia de profundidade $d$ e autenticação mútua, o voo de cada lado
cresce aproximadamente $d\,(p + s) + s$, com $p$ e $s$ os tamanhos de chave
pública e de assinatura. O acordo entre medida e contagem normativa é um segundo
indicador de validade interna, permite calcular o volume de combinações não
medidas e independe da plataforma, pois decorre de tamanhos normativos e da
estrutura do protocolo.

% ============================================================ DEGRAU
\section{Em Enlace Realista, o Custo é um Degrau}
\label{sec:degrau}

No perfil \texttt{lan} o volume praticamente não custa. P3 troca 32\% mais bytes
que H2 e leva o mesmo tempo, com diferença de 44~$\mu$s cujo intervalo contém o
zero, porque o gargalo é a CPU e não o canal. É esse o resultado que um enlace
ideal produz, e é por isso que ele induz otimismo.

Sob \texttt{wan} o quadro se inverte. C1 leva 85,5~ms, que são dois tempos de ida
e volta, um do estabelecimento TCP e outro do \textit{handshake} TLS~1.3, mais a
criptografia. H2 leva 92,6~ms, apenas 8\% acima, embora troque 8,56 vezes mais
bytes. Já P3 leva 157,7~ms, 84\% acima de C1. A decomposição incremental
inverte-se: o passo de H2 para P3, que sob \texttt{lan} respondia por 1,4\% do
acréscimo de tempo, responde por 90,1\% sob \texttt{wan}.

O crescimento não é proporcional aos bytes, é um degrau, e a última coluna da
Tabela~\ref{tab:resultados} indica onde ele está. A janela inicial de
congestionamento do Linux é de dez segmentos~\cite{rfc6928}; com MTU de 1500 e a
opção de \textit{timestamps} habilitada, que consome mais 12 bytes por segmento,
o tamanho máximo de segmento efetivo é 1\,448 bytes e a primeira janela comporta
14\,480. O voo do servidor mede 12\,160 bytes em H2 e 16\,042 em P3: H2 cabe, e
P3 excede o limiar por 1\,562 bytes, menos de 11\%. O critério, enunciado de
forma independente desta bancada, é que o custo extra aparece quando

\begin{equation}
d\,(p + s) + s \;>\; \mathit{initcwnd} \times \mathit{MSS}_{\mathrm{ef}} ,
\label{eq:limiar}
\end{equation}

\noindent
com o lado esquerdo dado pelo modelo da Seção~\ref{sec:bytes} e o direito pelos
parâmetros do enlace, sendo $\mathit{MSS}_{\mathrm{ef}}$ a MTU descontados os
cabeçalhos IP e TCP e as opções em uso. O que generaliza é o critério; o valor de
14\,480 bytes é específico da configuração medida.

Em autenticação mútua o critério aplica-se \emph{duas vezes}, porque cada lado
envia seu próprio voo. Em P3 o cliente envia 15\,877 bytes, que também excedem a
primeira janela, de modo que a configuração pode pagar até duas idas e voltas
adicionais. É consequência que a literatura de autenticação unilateral não teria
como observar, pois nela existe um único voo grande.

\begin{table}[!t]
\caption{Varredura da janela inicial sobre P3, perfil \texttt{wan}, três
repetições de 120 \textit{handshakes} por ponto. A previsão registrada antes da
coleta era de um degrau único entre dez e onze segmentos. Observaram-se dois
degraus, e o platô começa em doze.}
\label{tab:varredura}
\centering
\footnotesize
\begin{tabular}{@{}rrrr@{}}
\toprule
\textbf{\texttt{initcwnd}} & \textbf{janela} & \textbf{$T_{\mathrm{hs}}$ mediana}
 & \textbf{queda} \\
 & (bytes) & ($\mu$s) & ($\mu$s) \\
\midrule
10 & 14\,480 & 157\,363 & --- \\
11 & 15\,928 & 124\,858 & $-32\,506$ \\
12 & 17\,376 & \textbf{92\,490} & $-32\,368$ \\
13 & 18\,824 & 92\,686 & $+196$ \\
14 & 20\,272 & 92\,431 & $-255$ \\
20 & 28\,960 & 93\,692 & $+1\,261$ \\
40 & 57\,920 & 92\,213 & $-1\,479$ \\
\bottomrule
\end{tabular}
\end{table}

A Tabela~\ref{tab:varredura} leva o teste adiante, e o resultado é mais
informativo que a previsão. A resposta é de fato uma escada, e não um declive: de
doze segmentos em diante a mediana varia 2,0~ms em quinze medições, cerca de 2\%,
sem tendência. Mas há \emph{dois} degraus, de $32\,506$ e $32\,368$~$\mu$s,
iguais entre si a menos de 0,4\%, cada um da ordem de um tempo de ida e volta, e
o platô começa em doze e não em onze.

As duas discrepâncias têm a mesma origem e são verificáveis. A previsão usou
1\,460 bytes por segmento; com \textit{timestamps} habilitados, confirmados na
máquina, o valor efetivo é 1\,448. Os próprios degraus permitem estimá-lo sem
inspecionar a pilha: como o platô do voo do servidor, de 16\,042 bytes, só começa
em doze segmentos, tem-se $\mathit{MSS}_{\mathrm{ef}} < 16\,042/11 = 1\,458{,}4$;
e como o voo do cliente, de 15\,877 bytes, já cabe em onze, tem-se
$\mathit{MSS}_{\mathrm{ef}} \geq 15\,877/11 = 1\,443{,}4$. A varredura confina o
valor efetivo a uma faixa de quinze bytes que contém 1\,448 e exclui 1\,460. O
segundo degrau, por sua vez, é a manifestação direta da dupla aplicação do
critério: com dez segmentos ambos os voos excedem a janela e pagam-se duas idas e
voltas; com onze o voo do cliente passa a caber e paga-se uma; com doze cabem os
dois.

A intervenção, portanto, não apenas confirma o mecanismo como o quantifica: a
latência de P3 é função degrau do número de segmentos da primeira janela, com
patamares em duas, uma e zero idas e voltas adicionais, e a transição ocorre onde
a equação~\eqref{eq:limiar} a coloca quando avaliada com o tamanho de segmento
efetivo. A consequência é direta: em enlace realista o custo pós-quântico do
\textit{handshake} é quase nulo enquanto os voos couberem na primeira janela e
salta idas e voltas inteiras quando não couberem, de modo que trocar ML-DSA-65
por ML-DSA-44, reduzir a profundidade da cadeia ou elevar a janela inicial são
intervenções de efeito comparável, e a escolha entre elas é de engenharia de
implantação, não de criptografia.

% ============================================================ DISCUSSAO
\section{Discussão e Ameaças à Validade}
\label{sec:discussao}

O resultado é consistente com Paquin \textit{et
al.}~\cite{paquin2020benchmarking}, que apontavam o tamanho como fator dominante
sob condições adversas de rede, e acrescenta o mecanismo e o parâmetro que o
controla. A leitura prática é que o tamanho da cadeia importa menos por consumir
banda e mais por atravessar limiares do transporte, o que sugere que a
hibridização apenas na troca de chaves permanece o ponto de menor custo entre os
avaliados, já que seu acréscimo de 2\,270 bytes dificilmente cruza o limiar
sozinho.

Quanto à validade de conclusão e interna, as comparações vêm com intervalos por
\textit{bootstrap} hierárquico sobre execuções independentes, a ordem foi
randomizada e o contador de \textit{throttling} permaneceu em
zero~\cite{mytkowicz2009wrong}; o orçamento de bytes concorda com a contagem
normativa dentro de sete bytes, e a oscilação de comprimento aparece apenas nas
configurações ECDSA, como a codificação DER prevê.

A validade externa foi em parte delimitada na Seção~\ref{sec:plataforma}: as
magnitudes de tempo vêm de uma única microarquitetura e funcionam como limite
inferior do efeito, já que o termo criptográfico é ali inflado. Quanto ao degrau,
sua posição move-se com a janela inicial, com a MTU e com as opções TCP em uso,
de modo que ambientes que já elevam a janela, prática comum em provedores de
nuvem, podem não observar o efeito nas mesmas configurações; não depende, porém,
da disciplina de congestionamento, pelo argumento de partida lenta da
Seção~\ref{sec:config}. Defende-se como transportável a
equação~\eqref{eq:limiar} e o método de verificá-la, não o valor do limiar nem a
magnitude de 84\%. Quanto ao construto, $T_{\mathrm{hs}}$ mede o
\textit{handshake} a frio e não captura reuso de sessão, escolha deliberada
porque o regime de identidade efêmera por tarefa não o permite; RSA-3072 e
Falcon-512 não integram a campanha, e cadeias mistas, comuns em transição, não
foram avaliadas.

% ============================================================ CONCLUSAO
\section{Conclusões e Trabalhos Futuros}
\label{sec:conclusao}

Este artigo caracterizou o custo da migração pós-quântica no regime praticado por
pipelines DevSecOps, com autenticação mútua, conexões efêmeras e cadeias de dois
níveis, sobre plataforma nativa deliberadamente modesta. O orçamento de bytes é
aditivo, com resíduo abaixo de 0,03\%, e recuperável dos tamanhos normativos dos
esquemas, o que permite prever o volume de combinações não medidas e independe da
plataforma. E, em enlace com atraso e perda, a latência deixa de acompanhar o
tamanho e passa a ser um degrau no limiar da equação~\eqref{eq:limiar}: exceder a
janela inicial por menos de 11\% custou 84\% de latência. Em autenticação mútua o
critério incide sobre cada sentido, de modo que a resposta é uma escada de dois
degraus iguais, e a varredura estima o tamanho de segmento efetivo a partir da
latência, confinando-o a uma faixa de quinze bytes.

A implicação de projeto é que a escolha do esquema pós-quântico num pipeline não
deve ser feita apenas por comparação de custo criptográfico: calcular os voos
resultantes e confrontá-los com a janela inicial do ambiente alvo é um passo
barato, feito em papel, que pode tornar a escolha irrelevante diante de um ajuste
de transporte. E como o termo criptográfico encolhe com hardware mais capaz
enquanto o de rede não, a relevância relativa do degrau tende a crescer em
infraestrutura de maior porte.

Três direções decorrem das limitações: mapear o limiar variando MTU, inclusive em
enlaces com \textit{jumbo frames}, onde o critério prevê que o degrau desapareça;
replicar a campanha em nó de servidor contemporâneo, verificando a previsão de
que a razão entre os termos de rede e de criptografia cresce uma ordem de
grandeza; e medir uma variante baseada em KEMTLS~\cite{schwabe2020kemtls}, que
dispensa assinatura no \textit{handshake} e, pelo modelo apresentado, deveria
manter os voos abaixo do limiar nas condições em que os esquemas baseados em
assinatura o excedem.

\section*{Disponibilidade}
O código da bancada, os dados brutos por operação e os manifestos de configuração
do OpenSSL estão disponíveis em \texttt{FIXME:URL-DO-REPOSITORIO}.

% ============================================================ REFERENCIAS
\bibliographystyle{IEEEtran}
\bibliography{referencias}

\end{document}
